%% kapa: exam=1 round=27 year=2016 period=2 form=A subject=accounting subject_ko=회계학 seq=17 number=57
%% source: Q-Net 감정평가사 기출문제 (cst00309) artlSeq=5169402 file=extracted/1008/5169402/제27회 감정평가사 1차 문제/제27회 감정평가사 2교시문제(A형).pdf
%% conversion: kordoc 4.24.51 -> tools/gen1.py (auto); review=proofread-sonnet; flags=none
\begin{question}{57}
수익에 관한 설명으로 옳지 않은 것은?
\begin{choices}
  \item 용역제공거래의 성과를 신뢰성 있게 추정할 수 없는 경우에는 인식된 비용의 회수가능한 범위 내에서의 금액만을 수익으로 인식한다.
  \item 제품판매가격에 제품판매 후 제공할 용역에 대한 식별가능한 대가가 포함되어 있는 경우에는, 그 금액을 이연하여 용역수행기간에 걸쳐 수익으로 인식한다.
  \item 재화를 판매하고 동시에 당해 재화를 나중에 재구매하기로 하는 별도의 약정을 체결함으로써 판매거래의 실질적 효과가 상쇄되는 경우에는 두 개의 거래를 하나의 거래로 보아 회계처리한다.
  \item 판매대금의 회수가 구매자의 재판매에 의해 결정되는 경우에는 판매자가 소유에 따른 유의적인 위험을 부담하는 경우에 해당하므로, 당해 거래를 판매로 보지 아니하여 수익을 인식하지 아니한다.
  \item 수익은 기업이 받았거나 받을 경제적효익의 총유입을 의미하므로, 기업이 받는 판매세, 특정재화나 용역과 관련된 세금, 부가가치세 금액도 수익에 포함된다.
\end{choices}
\end{question}
\answer{5}
